\section{Probabilidade Condicional: de Reis a Prisioneiros}

\begin{quote}
The naked hulk alongside came,\\
And the twain were casting dice;\\
``The game is done! I've won! I've won!''\\
Quoth she, and whistles thrice.

\hfill Samuel Taylor Coleridge, \emph{The Rime of the Ancient Mariner}
\end{quote}

\subsection{Algumas regras de probabilidade. Probabilidade condicional}
\label{isaac:sec:prob-cond}

Esta parte segue o Capítulo 3 de \cite{Isaac1995}.

Vamos novamente lançar um par de dados uma vez e considerar o modelo probabilístico familiar, com seus 36 resultados possíveis e distribuição uniforme no espaço amostral resultante. Seja $A$ o evento ``obter 7''. No Capítulo 1 vimos que esse evento consiste em seis resultados, cada um com probabilidade $1/36$; portanto,
\[
\mathbb{P}(A)=\frac{6}{36}=\frac16.
\]
Seja $B$ o evento ``o primeiro dado (vermelho) mostra 1''. Listando os resultados de $B$, encontramos novamente seis resultados:
\[
(1,1),(1,2),(1,3),(1,4),(1,5),(1,6).
\]
Assim, $\mathbb{P}(B)=1/6$. E quanto a $\mathbb{P}(A\cup B)$? Por definição, $A\cup B$ ocorre quando obtemos 7, ou quando o primeiro dado mostra 1, ou quando ambas as coisas acontecem. Há 11 resultados: as seis maneiras de obter 7, entre as quais está $(1,6)$, mais os outros cinco resultados cujo primeiro componente é 1. Logo, $\mathbb{P}(A\cup B)=11/36$.

Já $A\cap B$ ocorre quando o primeiro dado mostra 1 e obtemos 7; isso só acontece quando aparece $(1,6)$. Portanto, $\mathbb{P}(A\cap B)=1/36$. Vale a fórmula
\begin{theorem}[Regra da adição]
\label{isaac:thm:adicao}
Para quaisquer eventos $A$ e $B$,
\[
\mathbb{P}(A\cup B)=\mathbb{P}(A)+\mathbb{P}(B)-\mathbb{P}(A\cap B).
\]
\end{theorem}

A razão é simples. Para calcular o lado esquerdo, somamos as probabilidades dos resultados que estão em $A$ ou em $B$, ou em ambos. No lado direito, $\mathbb{P}(A)+\mathbb{P}(B)$ conta duas vezes cada resultado que pertence simultaneamente a $A$ e $B$. Subtraímos então $\mathbb{P}(A\cap B)$ para que esses resultados sejam contados exatamente uma vez. A fórmula é completamente geral: vale em todo espaço de probabilidade e para todo par de eventos.

Se $A$ e $B$ não têm resultados em comum, isto é, $A\cap B=\varnothing$, então $\mathbb{P}(A\cap B)=0$ e
\[
\mathbb{P}(A\cup B)=\mathbb{P}(A)+\mathbb{P}(B).
\]
Nesse caso dizemos que os eventos são \emph{disjuntos}. Imagine $A$ e $B$ como dois terrenos que não se sobrepõem e as probabilidades como suas áreas: a área da união é a soma das áreas. Se os terrenos se sobrepõem, a propriedade total tem área menor que a soma das áreas individuais, e precisamos da fórmula mais geral.

Voltemos ao lançamento de dois dados, mas agora seja $A$ o evento ``obter soma 6''. Seus resultados são $(1,5),(2,4),(3,3),(4,2)$ e $(5,1)$; portanto, $\mathbb{P}(A)=5/36$. O evento $B$, ``obter 1 no primeiro dado'', ainda tem probabilidade $1/6$. Suponha agora que recebamos a informação de que o primeiro dado mostrou 1. Dada essa informação, qual é o valor de $\mathbb{P}(A)$?

É natural esperar que uma informação nova altere nossas ideias sobre as incertezas do experimento e nos faça reavaliar as probabilidades. Se sabemos que o primeiro dado mostrou 1, os únicos resultados possíveis são
\[
(1,1),(1,2),(1,3),(1,4),(1,5),(1,6).
\]
Podemos ignorar os demais, pois não podem ocorrer diante da informação recebida. Temos, assim, um novo espaço amostral; como partimos de uma distribuição uniforme, é natural supor que seus seis resultados continuam igualmente prováveis. Dado que o primeiro dado mostrou 1, o evento $A$ consiste no único resultado $(1,5)$ e parece ter probabilidade $1/6$. A informação adicional atualizou a probabilidade de $A$ de $5/36$ para $1/6$.

Em geral, quando chega uma informação nova, faz sentido atualizar o espaço de probabilidade e, com ele, o cálculo das probabilidades. As probabilidades atualizadas são chamadas \emph{probabilidades condicionais} dadas as novas informações.

\begin{definition}[Probabilidade condicional]
\label{isaac:def:prob-cond}
Se $\mathbb{P}(B)>0$, a probabilidade condicional de $A$ dado $B$ é
\[
\mathbb{P}(A\mid B)=\frac{\mathbb{P}(A\cap B)}{\mathbb{P}(B)}.
\]
\end{definition}

Se $B$ é conhecido, os resultados fora de $B$ não podem ocorrer; por isso restringimos nossa atenção a $A\cap B$. A definição torna essa quantidade proporcional a $\mathbb{P}(A\cap B)$. Como intuitivamente $\mathbb{P}(B\mid B)=1$, o fator de normalização correto é $\mathbb{P}(B)$. No exemplo acima,
\[
\mathbb{P}(A\mid B)=\frac{1/36}{6/36}=\frac16.
\]
Não definiremos probabilidades condicionais quando o evento condicionante tiver probabilidade zero. Multiplicando a fórmula por $\mathbb{P}(B)$, obtemos a forma útil
\[
\mathbb{P}(A\cap B)=\mathbb{P}(B)\,\mathbb{P}(A\mid B).
\]

A probabilidade condicional pode coincidir com a probabilidade original. Se $A$ é ``obter 7'' e $B$ é ``o primeiro dado mostra 1'', um cálculo simples dá $\mathbb{P}(A\mid B)=\mathbb{P}(A)=1/6$. Quando isso acontece, dizemos que $A$ é \emph{independente} de $B$: a informação fornecida por $B$ não trouxe informação nova sobre $A$. A independência será um dos grandes temas do Capítulo 5.

\subsection{O rei tem uma irmã?}
\label{isaac:sec:rei}

Considere o seguinte problema, para testar nossa habilidade com probabilidades condicionais (ele aparece como exercício em [28]):

\begin{example}[O rei e sua irmã]
\label{isaac:ex:rei}
O rei vem de uma família com dois filhos. Qual é a probabilidade de que o outro filho seja sua irmã?
\end{example}

Podemos tomar como espaço amostral o conjunto $S$ dos quatro pares $(B,B),(B,G),(G,B),(G,G)$, em que $B$ significa ``menino'', $G$ significa ``menina'', e as posições indicam, respectivamente, o primeiro e o segundo filho. Para resolver o problema precisamos fazer algumas suposições; mais uma vez, supomos que os quatro resultados são igualmente prováveis. Seja $U$ o evento ``um dos filhos é menina'' e $V$ o evento ``um dos filhos é o rei''. Queremos calcular $\mathbb{P}(U\mid V)$. Pela definição,
\[
\mathbb{P}(U\mid V)=\frac{\mathbb{P}(U\cap V)}{\mathbb{P}(V)}
=\frac{\mathbb{P}(\text{um filho é }B\text{ e o outro é }G)}{\mathbb{P}(V)}
=\frac{2/4}{3/4}=\frac23.
\]

Esse problema é traiçoeiro: muita gente pensa que a resposta deveria ser $1/2$, como no problema do carro e da cabra. Se a pergunta fosse ``qual é a probabilidade de que o irmão de uma pessoa seja uma irmã?'', a resposta seria $1/2$. Mas o enunciado informa sorrateiramente que um dos filhos, o rei, é homem; isso elimina $(G,G)$. Os três resultados restantes formam o espaço amostral condicional, e dois deles contêm um menino e uma menina. O problema ilustra mais uma vez a importância de interpretar cuidadosamente a informação transmitida pelo enunciado. Ambiguidades podem levar a espaços amostrais radicalmente diferentes e, portanto, a respostas diferentes.

\subsection{O dilema do prisioneiro}
\label{isaac:sec:prisioneiro}

Um bom exercício para pensar cuidadosamente sobre probabilidade condicional é o chamado \emph{dilema do prisioneiro}. Há também um problema famoso, completamente diferente e ligado à teoria dos jogos, com o mesmo nome. Considere três prisioneiros $A$, $B$ e $C$. Dois serão libertados, e os prisioneiros sabem disso, mas não sabem quais. O prisioneiro $A$ pede ao guarda que lhe diga o nome de um prisioneiro, diferente dele, que será libertado. O guarda se recusa e explica: ``Sua probabilidade de ser libertado agora é $2/3$. Se eu lhe disser que $B$, por exemplo, será libertado, você será um dos apenas dois prisioneiros cujo destino é desconhecido e sua probabilidade de libertação cairá para $1/2$. Como não quero prejudicar suas chances, não vou lhe dizer.'' O guarda está correto?

A resposta não é tão óbvia; é preciso investigar um pouco. O guarda pensa no espaço amostral das possibilidades de dois prisioneiros libertados:
\[
\{A,B\},\quad \{A,C\},\quad \{B,C\}.
\]
As chaves indicam pares não ordenados. Para obter um espaço de probabilidade, usamos a distribuição uniforme, supondo que o conselho de liberdade escolheu os prisioneiros ao acaso. Cada resultado tem probabilidade $1/3$, e a afirmação inicial do guarda, de que a probabilidade de libertação de $A$ é $2/3$, está correta. O problema começa quando ele parece afirmar que
\[
\mathbb{P}(A\text{ é libertado}\mid\text{o guarda diz que }B\text{ é libertado})=\frac12.
\]

O primeiro ponto é que essa probabilidade condicional não pode ser calculada no espaço amostral original: nele não existe o evento ``o guarda diz que $B$ será libertado''. Precisamos de um espaço mais complexo, que incorpore a fala do guarda. Considere os quatro resultados
\[
\begin{aligned}
O_1&=\{A,B,\text{o guarda diz que }B\text{ será libertado}\},\\
O_2&=\{A,C,\text{o guarda diz que }C\text{ será libertado}\},\\
O_3&=\{B,C,\text{o guarda diz que }B\text{ será libertado}\},\\
O_4&=\{B,C,\text{o guarda diz que }C\text{ será libertado}\}.
\end{aligned}
\]
Eles representam todas as possibilidades compatíveis com a libertação de dois prisioneiros e com uma afirmação possível do guarda. $O_1$ equivale a $\{A,B\}$, pois o guarda não tem escolha, e portanto $\mathbb{P}(O_1)=1/3$; analogamente, $\mathbb{P}(O_2)=1/3$.

Agora a situação fica interessante. A união $O_3\cup O_4$ é o evento $\{B,C\}$, de probabilidade $1/3$. Sem informação adicional, não há como determinar individualmente $\mathbb{P}(O_3)$ e $\mathbb{P}(O_4)$. Em geral, supõe-se que ambos sejam igualmente prováveis, com probabilidade $1/6$; isso corresponde ao guarda lançar uma moeda quando $B$ e $C$ são libertados. Mas ele poderia usar outro procedimento, por exemplo, sempre identificar $B$.

Primeiro suponhamos que cada um tenha probabilidade $1/6$. Então
\[
\begin{aligned}
\mathbb{P}(A\text{ é libertado}\mid\text{o guarda diz que }B\text{ é libertado})
&=\frac{\mathbb{P}(O_1)}{\mathbb{P}(\text{o guarda diz que }B\text{ é libertado})}\\
&=\frac{1/3}{1/3+1/6}=\frac23.
\end{aligned}
\]
Assim, a probabilidade condicional de libertação de $A$ é igual à probabilidade original. O termo $1/3+1/6=1/2$ aparece porque o evento ``o guarda diz que $B$ será libertado'' é $O_1\cup O_3$, uma união de eventos disjuntos. O mesmo argumento vale quando o guarda diz que $C$ será libertado, e novamente obtemos $2/3$. Portanto, a fala do guarda não alterou a probabilidade de libertação de $A$; os eventos ``$A$ é libertado'' e ``o guarda diz que $B$ será libertado'' são independentes.

Agora suponha que o guarda sempre identifique $B$ quando $B$ e $C$ forem libertados. Nesse caso, $\mathbb{P}(O_3)=1/3$ e $\mathbb{P}(O_4)=0$. O denominador passa a ser $1/3+1/3=2/3$, e a probabilidade condicional é $(1/3)/(2/3)=1/2$, exatamente como o guarda dissera. Em geral, ele pode identificar $B$ com qualquer probabilidade entre 0 e 1, e a probabilidade condicional ficará entre $1/2$ e 1.

Isso significa que o guarda controla o destino de $A$ pela maneira como formula sua fala? A conclusão contraria nossa intuição. Se ele sussurrasse a afirmação para si mesmo, em vez de contá-la a $A$, o mesmo raciocínio pareceria mostrar que poderia alterar o destino de $A$. A decisão foi tomada pelo conselho de liberdade, independentemente da fala do guarda. Isso sugere que devemos começar supondo a independência entre ``$A$ é libertado'' e a afirmação do guarda. Nesse caso, a probabilidade condicional é $2/3$, e a única possibilidade é atribuir probabilidade $1/6$ a cada um de $O_3$ e $O_4$. Essa é a solução que combina com o mundo real. As outras soluções são matematicamente corretas, mas não correspondem ao modelo adequado.

Observe que a probabilidade \emph{não condicional} de $A$ ser libertado é
\[
\mathbb{P}(A\text{ é libertado})=\mathbb{P}(O_1)+\mathbb{P}(O_2)=\frac23,
\]
independentemente da maneira como o guarda escolhe sua afirmação.

Uma versão do problema do carro e da cabra leva ao mesmo modelo matemático. Suponha que você sempre escolha inicialmente a porta 1, enquanto o carro e as cabras são distribuídos uniformemente atrás das três portas. O apresentador abre a porta 2 ou a porta 3, e você troca de porta. Qual é a probabilidade condicional de o carro estar atrás da porta 2, dado que o apresentador abriu a porta 3? Essa é a probabilidade de ganhar ao trocar.

Como mostrou Gillman [13], a resposta depende da probabilidade condicional de o apresentador abrir a porta 3 dado que o carro está atrás da porta 1. Identificando ``ganhar o carro'' com ``$A$ é libertado'' e ``abrir a porta 3'' com ``o guarda diz que $B$ é libertado'', os problemas são essencialmente equivalentes. A probabilidade condicional de ganhar varia entre $1/2$ e 1. A diferença é que, no dilema do prisioneiro, parece razoável supor desde o início que a libertação de $A$ é independente da fala do guarda. Já na nova versão do problema do carro e da cabra, a natureza dos eventos sugere dependência; por isso, qualquer valor entre $1/2$ e 1 pode ser uma resposta razoável.

\subsection{Tudo sobre urnas}
\label{isaac:sec:urnas}

Suponha que uma urna contenha dez bolas, seis vermelhas e quatro pretas, todas do mesmo tamanho e bem misturadas. Você escolhe uma bola ao acaso e anota sua cor. Não a repõe; depois de misturar novamente as bolas, escolhe uma segunda e anota sua cor. Defina
\[
A_1=\{\text{a primeira bola é vermelha}\},\qquad
A_2=\{\text{a segunda bola é vermelha}\}.
\]
Queremos as probabilidades de $A_1$, $A_1\cap A_2$ e $A_2$.

A probabilidade de $A_1$ é fácil. ``Bem misturadas'' é uma forma abreviada de dizer que supomos uma distribuição uniforme. O espaço amostral pode ser representado por dez resultados, um para cada bola, todos com probabilidade $0{,}1$. O evento $A_1$ contém seis resultados, logo $\mathbb{P}(A_1)=0{,}6$.

Para $A_1\cap A_2$, usamos a fórmula da probabilidade condicional:
\[
\mathbb{P}(A_1\cap A_2)=\mathbb{P}(A_1)\mathbb{P}(A_2\mid A_1)
=\frac6{10}\cdot\frac59=\frac13.
\]
O fator $5/9$ aparece porque, imediatamente antes da segunda escolha, restam cinco bolas vermelhas e quatro pretas, novamente bem misturadas.

Por fim, para calcular $\mathbb{P}(A_2)$ precisamos considerar todas as possibilidades da primeira escolha. O fato importante é
\[
A_2=(A_1\cap A_2)\cup(A_1^c\cap A_2),
\]
pois, se a segunda bola é vermelha, a primeira foi vermelha ou não foi vermelha, isto é, foi preta. Os eventos da união são disjuntos; portanto,
\[
\mathbb{P}(A_2)=\mathbb{P}(A_1\cap A_2)+\mathbb{P}(A_1^c\cap A_2).
\]
O primeiro termo é $1/3$. Para o segundo,
\[
\mathbb{P}(A_1^c\cap A_2)=\mathbb{P}(A_1^c)\mathbb{P}(A_2\mid A_1^c)
=\frac4{10}\cdot\frac69=\frac4{15}.
\]
Assim, $\mathbb{P}(A_2)=1/3+4/15=0{,}6$.

É interessante que calculamos a probabilidade de uma interseção usando uma probabilidade condicional dada. Nos problemas anteriores, calculávamos uma probabilidade condicional a partir de probabilidades de interseções. Aqui a probabilidade que surge naturalmente é condicional: a probabilidade da segunda escolha depende do que ocorreu na primeira.

Outro fato interessante é que $A_1$ e $A_2$ têm a mesma probabilidade, $0{,}6$. Isso não é coincidência. Se há $a$ bolas vermelhas e $b$ pretas, com pelo menos duas bolas, então cada seleção tem probabilidade
\[
\mathbb{P}(\text{bola vermelha na seleção }n)=\frac{a}{a+b},
\]
desde que existam bolas suficientes para realizar $n$ seleções.

À primeira vista, isso parece curioso: a segunda seleção ocorre depois da primeira, e o tempo esconde uma simetria importante. Em vez de pensar nas escolhas em sequência, imagine que escolhemos uma única vez, alcançando simultaneamente a urna com as duas mãos e retirando uma bola com cada mão. O resultado pode ser um par ordenado como $(R,B)$, em que a primeira coordenada é a cor da bola na mão esquerda e a segunda é a cor da bola na mão direita. A mão esquerda pode ser identificada com a primeira bola e a direita com a segunda. A distribuição dos pares é uniforme, como nos modelos do Capítulo 2. Todo par com $R$ na primeira coordenada corresponde, ao trocar as coordenadas, a um par com $R$ na segunda. Logo, $A_1$ e $A_2$ têm o mesmo número de resultados e a mesma probabilidade.

O mesmo raciocínio mostra por que esperamos o mesmo fenômeno na $n$-ésima seleção. Essa maneira de olhar para o problema revela ideias que não aparecem tão facilmente na descrição sequencial.

O modelo de seleção sem reposição pode ser modificado. Por exemplo, podemos escolher uma bola e, se ela for vermelha, recolocá-la e acrescentar outra bola vermelha; se for preta, recolocá-la e acrescentar outra preta. O número total de bolas cresce, em vez de diminuir. Esse é um caso particular do chamado esquema de urnas de Pólya, um modelo rudimentar de doença contagiosa. Cada seleção representa a amostragem de um indivíduo; a bola vermelha representa uma pessoa infectada e a preta, uma pessoa saudável. Cada infecção descoberta aumenta a probabilidade de encontrar outra pessoa infectada, enquanto cada pessoa saudável descoberta aumenta a probabilidade de encontrar outra saudável. Versões refinadas desse modelo permitem estudar a evolução da doença a longo prazo.

\subsection{Exercícios}
\label{isaac:sec:exercicios}

\begin{exercise}
\label{isaac:ex:dados-eventos}
Seja $S$ o espaço de probabilidade usual de um par de dados lançado uma vez, com distribuição uniforme. Suponha que $A$ seja o evento ``o primeiro dado é ímpar'' e $B$ o evento ``o segundo dado é par''. Descreva em palavras cada evento a seguir e calcule sua probabilidade: (a) $A\cap B$; (b) $(A\cap B)\cup(A^c\cap B^c)$; (c) $A^c$; (d) $(A\cup B)^c$.
\end{exercise}

\begin{exercise}
\label{isaac:ex:soma-11}
Lance um par de dados uma vez. Qual é a probabilidade de obter 11? Qual é a probabilidade de obter 11 dado que a soma das faces é ímpar? E qual é a probabilidade de obter 11 dado que a soma das faces é um número ímpar maior que 3?
\end{exercise}

\begin{exercise}
\label{isaac:ex:moeda-quatro}
Lance uma moeda quatro vezes. Determine a probabilidade de obter pelo menos duas caras. Determine também a probabilidade de obter pelo menos duas caras dado que ocorreu pelo menos uma cara, e a probabilidade de obter caras nos quatro lançamentos dado que ocorreram pelo menos duas caras.
\end{exercise}

\begin{exercise}
\label{isaac:ex:urna-reposicao}
Uma urna contém cinco bolas vermelhas e cinco pretas. Uma bola é escolhida ao acaso e, juntamente com outra bola da mesma cor, é devolvida à urna. Uma segunda bola é então escolhida ao acaso. Calcule as probabilidades de que (a) a primeira bola seja vermelha e a segunda vermelha; (b) a primeira bola seja vermelha e a segunda preta; (c) a segunda bola seja vermelha; (d) a segunda bola seja preta.
\end{exercise}

\begin{exercise}
\label{isaac:ex:regra-produto}
Sejam $A$, $B$ e $C$ eventos tais que $A$ e $A\cap B$ tenham probabilidade positiva. Usando a fórmula da probabilidade condicional, prove que
\[
\mathbb{P}(A)\mathbb{P}(B\mid A)\mathbb{P}(C\mid A\cap B)
=\mathbb{P}(A\cap B\cap C).
\]
\end{exercise}
